\documentclass[aps,pra,reprint,amsmath,amssymb,nofootinbib,longbibliography,floatfix]{revtex4-2}
\usepackage{graphicx}
\usepackage{bm}
\usepackage[colorlinks=true,citecolor=blue,linkcolor=blue,urlcolor=blue]{hyperref}
\newcommand{\ket}[1]{\lvert #1\rangle}
\newcommand{\bra}[1]{\langle #1\rvert}
\newcommand{\Tr}{\operatorname{Tr}}
\newcommand{\I}{\mathbb{I}}
\newcommand{\thv}{\bm{\theta}}
\newcommand{\est}{\hat{\bm{\theta}}}
\begin{document}
\title{Exact single-use Bayesian Nagaoka--Hayashi attainment for noisy qubit rotations}
\author{Xiaojian Zeng}
\email{xiaojianzeng@foxmail.com}
\affiliation{PBC School of Finance, Tsinghua University, Beijing 100190, China}
\date{October 12, 2026}
\begin{abstract}
We determine the exact Bayesian precision for simultaneous estimation of
three qubit-rotation coordinates using one fixed Bell probe and a noiseless
ancilla. The rotation is followed by known depolarization of strength
$0\leq\lambda\leq1$, the prior is uniform Cartesian volume on the
radius-$\pi/4$ parameter ball, and the loss is one third of the squared
coordinate error. The optimum over all joint measurements and estimators
equals the Bayesian Nagaoka--Hayashi relaxation and is
$\pi^2/80-4(1-\lambda)^2[3/(2\pi)-12/\pi^3]^2$.
Four orthogonal maximally entangled outcomes attain this value: a fixed
local rotation followed by a Bell measurement suffices, while the posterior
estimates shrink with noise. A covariant sphere measurement is also optimal.
We give algebraic positive-semidefinite certificates for both the physical
risk and the relaxation, valid throughout the noise interval, including
its pure-state and no-information endpoints. The result converts the
previous single-use numerical near-attainment into a constructive exact
precision limit. It concerns the fixed single-use experiment; attainment
for two through four parallel uses is not established.
\end{abstract}
\maketitle

\section{Introduction}
Prior information and measurement incompatibility both matter when only
a finite number of quantum observations are available. For quadratic
Bayesian loss, prior moments determine the optimal posterior estimator
and a symmetric-logarithmic-derivative (SLD) lower bound
\cite{Rubio2020}. The Nagaoka--Hayashi (NH) semidefinite formulation
\cite{Conlon2021} has a Bayesian extension \cite{Suzuki2024} that can be
tighter, but the relaxed measurement moments need not by themselves
specify a physical measurement. Recent work relates this distinction to
Bayesian measurement incompatibility \cite{Albarelli2026}.

Bia and Gagatsos studied simultaneous estimation of all three qubit-rotation
coordinates with a uniform ball prior and depolarizing noise
\cite{Bia2026}. For up to four parallel uses, their computed strategies
nearly attain the Bayesian NH bound, both with optimized probes and with
fixed tensor products of Bell probes. They explicitly characterize this
as numerical evidence rather than an attainment theorem. We settle the
fixed-Bell, single-use sector exactly for every allowed noise strength.
The contribution is a finite attaining measurement and matching algebraic
certificates, together with a closed precision formula.

The scalar-plus-triplet prior moments used below have a close counterpart
in the phase-imaging example of Ref.~\cite{Albarelli2026}, whose NH
attainment is also assessed numerically. This representation and the
general Bayesian optimality framework are established tools. Likewise,
NH-saturating noisy-rotation measurements in a local, unbiased,
phase-damping model were constructed in Ref.~\cite{Conlon2021}.
Our result concerns a finite volume prior, depolarization, and posterior
estimation without an unbiasedness condition. We do not extend it to
multiple uses, other probes, other losses, or adaptive experiments.

\section{Experiment and exact prior moments}
The system qubit and an untouched ancillary qubit are initially in
$\ket{\Phi^+}=(\ket{00}+\ket{11})/\sqrt2$. The system undergoes
\begin{equation}
 U_{\thv}=e^{-i\thv\cdot\bm{\sigma}},\qquad
 \mathcal D_\lambda(\rho)=(1-\lambda)\rho+\lambda\I_2/2,
 \label{eq:channel}
\end{equation}
with known $\lambda\in[0,1]$. The Pauli eigenvalues are $\pm1$;
there is no factor $1/2$ in the exponent. The Bloch-vector rotation
angle is consequently $2\|\thv\|$. In system--ancilla order the
normalized output is
\begin{align}
 \rho_{\lambda,\thv}&=(1-\lambda)\ket{v_{\thv}}\bra{v_{\thv}}
                  +\lambda\I_4/4,\label{eq:state}\\
 \ket{v_{\thv}}&=(U_{\thv}\otimes\I_2)\ket{\Phi^+}.
\end{align}
This is a trace-one state from a completely positive trace-preserving
channel. In the input--output Choi order of Ref.~\cite{Bia2026}, if $W$
swaps the factors, $\rho_{\lambda,\thv}=WJ_{\thv}W^\dagger/2$;
the unnormalized one-use Choi operator has trace two.

Put $a=\pi/4$. The prior and Bayes risk are
\begin{align}
 p(\thv)&=\frac{3}{4\pi a^3}\,\mathbf1_{\{\|\thv\|\leq a\}},
                  \label{eq:prior}\\
 R(M,\est)&=\frac13\mathbb E\bigl[\|\est-\thv\|^2\bigr].
                  \label{eq:risk}
\end{align}
The expectation includes the prior and measurement outcomes. We allow
arbitrary joint positive operator-valued measures (POVMs), finite or
continuous, and measurable estimators with finite second moment.
All coordinates and costs are dimensionless. The prior is uniform
Cartesian volume, with radial density $3r^2/a^3$; it is neither Haar
measure nor a uniform radial density. Its no-information risk is
\begin{equation}
 R_{\mathrm p}=\frac{\mathbb E\|\thv\|^2}{3}
             =\frac{a^2}{5}=\frac{\pi^2}{80}.
 \label{eq:priorrisk}
\end{equation}

For a transparent real-coordinate representation define
\begin{equation}
 \ket0=\ket{\Phi^+},\qquad
 \ket{\widetilde i}=-i(\sigma_i\otimes\I_2)\ket{\Phi^+}
 \quad(i=1,2,3).
 \label{eq:basis}
\end{equation}
These four physical vectors are orthonormal. For $\thv=r\bm n$,
\begin{equation}
 \ket{v_{\thv}}=\cos r\ket0+\sin r\sum_i n_i\ket{\widetilde i}.
 \label{eq:realstate}
\end{equation}
The rephasing in Eq.~\eqref{eq:basis} changes the basis, not the sign
of the estimated parameter. Write $E_{ab}=\ket a\bra b$ in this basis,
with indices $a,b=0,1,2,3$, and define
\begin{equation}
 P_0=E_{00},\quad P_1=\sum_{i=1}^3E_{ii},\quad F_i=E_{0i}+E_{i0}.
\end{equation}
The angular identities $\mathbb E[n_i]=0$ and
$\mathbb E[n_i n_j]=\delta_{ij}/3$ give
\begin{align}
 \Gamma_0&=\int p(\thv)\rho_{\lambda,\thv}\,d^3\thv
             =A P_0+B P_1,\label{eq:moments}\\
 \Gamma_i&=\int p(\thv)\theta_i\rho_{\lambda,\thv}\,d^3\thv
             =QF_i,\nonumber\\
 A&=(1-\lambda)C+\lambda/4,\nonumber\\
 B&=(1-\lambda)(1-C)/3+\lambda/4,\nonumber\\
 Q&=(1-\lambda)q,\qquad A+3B=1.\label{eq:ABQ}
\end{align}
The radial integrals evaluate exactly to
\begin{align}
 C&=\frac{3}{a^3}\int_0^a r^2\cos^2r\,dr
       =\frac12+\frac3\pi-\frac{24}{\pi^3},\label{eq:C}\\
 q&=\frac{1}{a^3}\int_0^a r^3\sin r\cos r\,dr
       =\frac{3}{2\pi}-\frac{12}{\pi^3}.\label{eq:q}
\end{align}
Both $A$ and $B$ are strictly positive on the entire noise interval.
At zero noise the individual states are pure, while their prior average
is full rank.

\section{A finite attaining measurement}
Let the four sign vectors be
\begin{equation}
 \begin{aligned}
 \bm s_0&=(1,1,1),&\bm s_1&=(1,-1,-1),\\
 \bm s_2&=(-1,1,-1),&\bm s_3&=(-1,-1,1).
 \end{aligned}
\end{equation}
Define the outcome vectors and effects
\begin{equation}
 \ket{e_j}=\frac{\ket0+\sum_i s_{ji}\ket{\widetilde i}}{2},\qquad
 M_j=\ket{e_j}\bra{e_j}.
 \label{eq:pvm}
\end{equation}
Since $1+\bm s_j\cdot\bm s_l=4\delta_{jl}$, this is an orthonormal
projective measurement on the full two-qubit output space. The outcome
probabilities and first moments are
\begin{equation}
 p_j=\Tr(\Gamma_0M_j)=\frac14,\qquad
 g_{ij}=\Tr(\Gamma_iM_j)=\frac Q2s_{ji}.
\end{equation}
For quadratic loss the optimal estimator at fixed measurement is the
posterior mean, so
\begin{equation}
 \est_j=\frac{\bm g_j}{p_j}=2Q\bm s_j.
 \label{eq:est}
\end{equation}
The posterior decomposition of squared error yields
\begin{align}
 R(M,\est)&=R_{\mathrm p}-\frac13\sum_jp_j\|\est_j\|^2
                =R_{\mathrm p}-4Q^2.\label{eq:achieved}
\end{align}
The measurement does not depend on $\lambda$; only the estimates do.
They remain in the prior ball: from
$q=\mathbb E[r\sin r\cos r]/3\leq\mathbb E[r]/6=a/8$,
\begin{equation}
 \|\est_j\|=2\sqrt3Q\leq\frac{\sqrt3a}{4}<a.
\end{equation}

The measurement also has a simple two-qubit implementation. Set
\begin{equation}
 V_j=\frac{\I_2-i\bm s_j\cdot\bm\sigma}{2},\qquad
 \ket{e_j}=(V_j\otimes\I_2)\ket{\Phi^+}.
 \label{eq:V}
\end{equation}
Each $V_j$ is unitary. For the stated ordering,
\begin{equation}
 V_0^\dagger V_j=\I_2,\ i\sigma_y,\ i\sigma_z,\ i\sigma_x
 \quad (j=0,1,2,3).
 \label{eq:circuit}
\end{equation}
Apply $V_0^\dagger$ to the system after the noisy channel and then perform
a Bell-basis measurement, with the label permutation in
Eq.~\eqref{eq:circuit}. This implements Eq.~\eqref{eq:pvm}.
It is a circuit identity requiring a joint Bell measurement, rather
than a claim of an experimental implementation or a product measurement.

\section{Exact optimality and NH saturation}
OpenAI Codex assisted with the derivations and verification code used here
and in the appendices. Human scientific verification and final tool-version
disclosure remain pending for this local draft.

For the fixed probe, the Bayesian NH program of Refs.~\cite{Suzuki2024,Bia2026}
is the minimization of
\begin{equation}
 \mathcal R(L,X)=R_{\mathrm p}
 +\frac13\sum_i\Tr(\Gamma_0 L_{ii})
 -\frac23\sum_i\Tr(\Gamma_i X_i)
 \label{eq:nhobj}
\end{equation}
over Hermitian $X_i$ and blocks $L_{ij}=L_{ji}=L_{ij}^\dagger$, subject to
\begin{equation}
 S(L,X)=\begin{pmatrix}L&X\\X^\dagger&\I_4\end{pmatrix}\succeq0,
 \qquad X=\begin{pmatrix}X_1\\X_2\\X_3\end{pmatrix}.
 \label{eq:schur}
\end{equation}
Every physical measurement with real estimates supplies feasible moments
$X_i=\int t_i\,dM$ and $L_{ij}=\int t_it_j\,dM$:
\begin{equation}
 S(L,X)=\int
 \begin{pmatrix}t_1\I_4\\t_2\I_4\\t_3\I_4\\\I_4\end{pmatrix}
 dM
 \begin{pmatrix}t_1\I_4&t_2\I_4&t_3\I_4&\I_4\end{pmatrix}.
 \label{eq:physicalgram}
\end{equation}
Thus its optimum $C_{\mathrm{NH}}$ is a lower bound on the physical
optimum $R_1$. To prove tightness, we give a dual matrix independently
of a numerical solution.

In $\mathbb C^3\otimes\mathbb C^4$, the first factor labels parameters
and the second uses Eq.~\eqref{eq:basis}. Define
\begin{align}
 \ket{\omega}&=\sum_{i=1}^3\ket i\otimes\ket{\widetilde i},\nonumber\\
 \ket{h_{ij}}&=\ket i\otimes\ket{\widetilde j}
                   -\ket j\otimes\ket{\widetilde i},\nonumber\\
 K&=\frac A3\sum_i\ket{i,0}\bra{i,0}
       +\frac B3\ket\omega\bra\omega
       +\frac B3\sum_{i<j}\ket{h_{ij}}\bra{h_{ij}}.
 \label{eq:Kgram}
\end{align}
This is positive semidefinite. Its $4\times4$ Hilbert-space blocks are
\begin{equation}
 K_{ij}=\frac A3\delta_{ij}P_0+
        \frac B3\bigl(\delta_{ij}P_1+E_{ij}-E_{ji}\bigr).
 \label{eq:Kblocks}
\end{equation}
Let $x=2Q$ and let $X^\star$ be the column of the three operators
$X_i^\star=xF_i$. These operators define the certificate; they need
not equal the finite measurement's first moments. The lifted matrix
\begin{align}
 Z&=\begin{pmatrix}\I_{12}\\-X^{\star\dagger}\end{pmatrix}
       K\begin{pmatrix}\I_{12}&-X^\star\end{pmatrix}\succeq0,
       \label{eq:Z}\\
 (KX^\star)_i&=x\left(\frac A3E_{0i}+BE_{i0}\right),\nonumber\\
 Z_{44}&=x^2\left(3BP_0+\frac A3P_1\right)\label{eq:Z44}
\end{align}
obeys, by $A+3B=1$,
\begin{align}
 Z_{ii}&=\Gamma_0/3,\nonumber\\
 Z_{ij}+Z_{ji}&=0\quad(i<j\leq3),\nonumber\\
 Z_{i4}+Z_{4i}&=-2\Gamma_i/3,\nonumber\\
 \Tr Z_{44}&=4Q^2.\label{eq:dualaffine}
\end{align}
For any feasible $S(L,X)$, these equations give the exact weak-duality
identity
\begin{equation}
 \mathcal R(L,X)-(R_{\mathrm p}-4Q^2)
             =\Tr[ZS(L,X)]\geq0.
 \label{eq:weakdual}
\end{equation}
The off-diagonal $K$ contributions cancel because $L_{ij}=L_{ji}$;
the Hermitian $X_i$ give precisely the first-moment coefficients in
Eq.~\eqref{eq:nhobj}. Consequently
\begin{equation}
 R_{\mathrm p}-4Q^2\leq C_{\mathrm{NH}}\leq R_1
          \leq R(M,\est)=R_{\mathrm p}-4Q^2.
\end{equation}
All inequalities are equalities, proving the stated result:
\begin{equation}
 \boxed{\begin{aligned}
 R_1(\lambda)&=C_{\mathrm{NH}}(1,\lambda)\\
 &=\frac{\pi^2}{80}-4(1-\lambda)^2
       \left(\frac{3}{2\pi}-\frac{12}{\pi^3}\right)^2.
 \end{aligned}}
 \label{eq:result}
\end{equation}
Positivity follows from the Gram factors in Eqs.~\eqref{eq:Kgram}
and \eqref{eq:Z}, so no solver tolerance or noise grid enters the proof.
A separate physical loss-operator certificate, valid for every real
estimate, is given in Appendix~\ref{app:physical}.

\begin{figure}[htbp]
 \includegraphics[width=\columnwidth]{manuscript/figures/risk.pdf}
 \caption{Dimensionless mean-square coordinate risk for one fixed Bell
 probe and the radius-$\pi/4$ volume prior. The optimal risk and NH bound
 are the same solid curve, Eq.~\eqref{eq:result}; the dashed curve is
 the Bayesian SLD bound, Eq.~\eqref{eq:sld}; the dotted line is the
 prior risk. The shaded separation is the exact incompatibility penalty.
 The curves are evaluations of analytic formulas, not measurement data
 or numerical optimization results. OpenAI Codex assisted with the
 plotting code; the figure was rendered from the accompanying formula table.}
 \label{fig:risk}
\end{figure}

\section{Covariance and the incompatibility penalty}
An alternative optimal measurement is continuous. With the normalized
sphere measure $d\mu(\bm n)=d\Omega/(4\pi)$, set
\begin{align}
 \ket{m(\bm n)}&=\ket0+\sqrt3\sum_i n_i\ket{\widetilde i},\nonumber\\
 dM(\bm n)&=\ket{m(\bm n)}\bra{m(\bm n)}\,d\mu(\bm n).
 \label{eq:sphere}
\end{align}
Spherical first and second moments imply $\int dM=\I_4$.
The outcome probability density relative to $d\mu$ is $A+3B=1$,
and the posterior mean is $2\sqrt3Q\bm n$. Its risk is therefore
again $R_{\mathrm p}-4Q^2$. This POVM is covariant under the scalar
and vector representation of $\mathrm{SO}(3)$. The four-outcome
measurement is a finite optimal alternative; it does not discretize
the continuous parameter prior.

The finite measurement's moment operators make its mechanism explicit.
For $\{i,j,l\}=\{1,2,3\}$,
\begin{equation}
 X_i^{\rm tet}=2Q(F_i+E_{jl}+E_{lj}),\qquad
 (X_i^{\rm tet})^2=4Q^2\I_4.
 \label{eq:completion}
\end{equation}
They commute pairwise. The triplet completion terms have zero overlap
with the signal first moments $\Gamma_i$ and enable a common projective
measurement. They do not turn the SLD operators into commuting operators.
Indeed the Bayesian Lyapunov equations
$(\Gamma_0S_i+S_i\Gamma_0)/2=\Gamma_i$ give
\begin{align}
 S_i&=\frac{2Q}{A+B}F_i,\nonumber\\
 C_{\mathrm{SLD}}&=R_{\mathrm p}-\frac{4Q^2}{A+B},\nonumber\\
 R_1-C_{\mathrm{SLD}}&=\frac{8BQ^2}{A+B}>0\quad(\lambda<1).
 \label{eq:sld}
\end{align}
This is a finite Bayesian incompatibility cost, rather than an asymptotic
or locally unbiased variance. Figure~\ref{fig:risk} displays these exact
expressions.



At zero noise, the formulas give
\begin{align}
 R_1(0)&\simeq0.09064783829,\nonumber\\
 C_{\mathrm{SLD}}(0)&\simeq0.08180546613,\nonumber\\
 R_1(0)/C_{\mathrm{SLD}}(0)&\simeq1.108090236.\label{eq:benchmark}
\end{align}
This reproduces the rounded single-use ratio $1.108$ reported in
Ref.~\cite{Bia2026}. For one use, rotation averaging of a tester preserves
the isotropic prior and loss, and irreducibility of the qubit input
representation fixes its invariant marginal to $\I_2/2$; hence the
source's invariant single-use benchmark maps to this Bell experiment.
This benchmark comparison does not assert optimality of Bell probes
for multiple uses.

\section{Scope and physical interpretation}
The optimal apparatus can be fixed while the known depolarization changes:
the posterior vectors are multiplied by $1-\lambda$, and the precision
gain below the prior risk scales as $(1-\lambda)^2$. At $\lambda=0$ the
proof applies directly to pure individual outputs, using positive prior
moments and exact Gram factors. At $\lambda=1$, $Q=0$ and all estimates
are zero; the physical optimum, NH bound, and SLD bound are the prior
risk. No full-rank endpoint continuation, perturbative expansion, or
large-sample approximation is needed.

The closest phase-imaging moments in Ref.~\cite{Albarelli2026},
Eqs.~(58)--(59), have the same star structure after rephasing: for
three parameters, $\alpha^2=A/B$ and a first-moment rescaling match
the measurement optimization, while the prior-loss constants differ.
Thus the star reduction is not an additional contribution of this work.
The finite measurement and positive NH certificate supply the exact
attainment step at the present rotation-model scope.

Although the source numerical study extends through four parallel uses,
this paper concerns one use with a fixed Bell probe. In a multiple-use
experiment every output encodes the same unknown vector, so
$\int p(\thv)\rho_{\lambda,\thv}^{\otimes k}d^3\thv$ is generally
not $\Gamma_0^{\otimes k}$. The present four-dimensional certificates
therefore do not settle NH attainment for $k=2,3,4$. They also do not
establish a global optimum over arbitrary probes or adaptive protocols.

\section*{Data availability}
The analytic expressions and optimality certificates are given in this
article. Symbolic checks, direct measurement-replay code, and the code
and tabulated formula evaluations for Fig.~\ref{fig:risk} accompany
the local draft. No experimental data were used.

\section*{Acknowledgments and use of artificial intelligence}
OpenAI Codex assisted with literature checks, analytic derivations,
symbolic and numerical verification code, figure generation, and
manuscript preparation. Automated checks included a separately
reconstructed calculation in a fresh AI context. Human scientific
review and approval of this local draft, including confirmation of
tool versions and the final disclosure, remain pending.

\appendix
\section{A physical loss-operator certificate}\label{app:physical}
Define the second-moment loss operator
\begin{equation}
 H=\frac13\int p(\thv)\|\thv\|^2\rho_{\lambda,\thv}\,d^3\thv,
 \qquad\Tr H=R_{\mathrm p}.
\end{equation}
For a reported vector $\bm t\in\mathbb R^3$, its loss operator is
\begin{equation}
 \mathcal F(\bm t)=H+\frac13\left(\|\bm t\|^2\Gamma_0
                                 -2\sum_it_i\Gamma_i\right).
\end{equation}
The general supporting-operator criterion is part of Bayesian quantum
decision theory; see Appendix~B of Ref.~\cite{Albarelli2026}. Here a
specific supporting operator is
\begin{equation}
 Y=H-12BQ^2P_0-\frac43AQ^2P_1.
\end{equation}
Rotational covariance lets us take $\bm t=u\bm e_3$, where
$u=\|\bm t\|\geq0$. The $0,3$ block of $\mathcal F(\bm t)-Y$ is
\begin{equation}
 \frac13\begin{pmatrix}
 Au^2+36BQ^2&-2Qu\\-2Qu&Bu^2+4AQ^2
 \end{pmatrix};
\end{equation}
the two transverse entries are $(Bu^2+4AQ^2)/3$. All diagonal entries
are nonnegative, and the block determinant is
\begin{equation}
 \det[\mathcal F(\bm t)-Y]_{0,3}
             =\frac{AB}{9}(u^2-12Q^2)^2\geq0.
\end{equation}
Thus $\mathcal F(\bm t)\succeq Y$ for every real estimate, including
estimates outside the prior ball. Integrating against any POVM proves
$R\geq\Tr Y=R_{\mathrm p}-4Q^2$. At $u=2\sqrt3Q$ the block annihilates
$(1,\sqrt3)^T$, giving equality for the sphere measurement and the
four tetrahedral directions. This certificate proves physical
optimality without assuming that every feasible NH moment pair is
realizable.

\section{Explicit covariant NH primal}\label{app:primal}
For completeness, the sphere measurement supplies
\begin{align}
 X_i^{\rm sph}&=2QF_i,\nonumber\\
 L_{ij}^{\rm sph}&=4Q^2\delta_{ij}P_0
       +\frac{12Q^2}{5}(\delta_{ij}P_1+E_{ij}+E_{ji}).
\end{align}
These formulas follow from
$\mathbb E[n_in_jn_an_b]=(\delta_{ij}\delta_{ab}+
\delta_{ia}\delta_{jb}+\delta_{ib}\delta_{ja})/15$.
The Schur complement $L^{\rm sph}-X^{\rm sph}X^{\rm sph\dagger}$
vanishes on the scalar Hilbert sector. On the parameter--triplet
tensor space it vanishes on trace and antisymmetric tensors and
equals $24Q^2/5$ times the identity on symmetric traceless tensors.
This gives another direct positive-semidefinite primal with objective
$R_{\mathrm p}-4Q^2$, including $Q=0$.

\bibliographystyle{apsrev4-2}
\bibliography{manuscript/references}
\end{document}
